% Part:sets-functions-relations
% Chapter: size-of-sets
% Section: introduction

\documentclass[../../../include/open-logic-section]{subfiles}

\begin{document}

\olfileid{sfr}{siz}{int}

\olsection{تعارف}

جدوں گیورگ کانتور نے 1870واں دے دہاکے وچ سیٹ نظریہ تیار کیتا، اوہدے
مقصداں وچوں اک لامتناہی اکٹھ---جیویں وچلے زمانے دے فلسفی آکھدے سن،
اک بالفعل لامتناہیت---دے خیال نوں قابلِ قبول بنانا سی۔ ایس کم دا اک
اہم حصہ وکھ وکھ سیٹاں دے \emph{سائز} نال اوہدا ورتاؤ سی۔ جے $a$،
$b$ تے $c$ سارے اک دوجے توں وکھ نیں، تاں سیٹ $\{a, b, c\}$ بدیہی
طور اُتے $\{a, b\}$ نالوں \emph{وڈا} اے۔ پر لامتناہی سیٹاں دا کیہ؟
کی اوہ سارے اک دوجے جڈے وڈے نیں؟ پتا لگدا اے کہ نہیں۔

ایتھے پہلا اہم خیال گنتی وار فہرست دا اے۔ اسی ہر متناہی سیٹ دے سارے
!!{element}s نوں وار وار لکھ کے اوہدی فہرست بنا سکدے آں۔ کجھ لامتناہی
سیٹاں دے وی سارے !!{element}s دی فہرست بنا سکدے آں، جے فہرست نوں آپ
لامتناہی ہون دی اجازت ہووے۔ ایہو جہے سیٹ، یعنی گنتی وار فہرست بنن
جوگ سیٹ، !!{enumerable} آکھلاندے نیں۔ کانتور دا حیران کرن والا نتیجہ،
جیہنوں اسی ایس باب دے آخر تک پوری طرح سمجھ لاواں گے، ایہ سی کہ کجھ
لامتناہی سیٹ !!{enumerable} نہیں ہوندے۔

\end{document}
